\MathBookSourcePage{278}
\subsection{习题}
\label{sec:ch09-exercises}
\setcounter{exerciseitem}{0}

\begin{MBExercise}
\label{exr:ch09-01}
证明在二维或更高维中, 非退化网格是局部拟一致的. 提示: 相邻单元都可以通过一列具有公共面的单元彼此连接.
\end{MBExercise}

\begin{MBExercise}
\label{exr:ch09-02}
证明定理~\ref{thm:ch09-mesh-gradient-l2-error}.
\end{MBExercise}

\begin{MBExercise}
\label{exr:ch09-03}
对图~\ref{fig:ch09-highly-graded-meshes} 中的网格计算 $\|\nabla h\|_{L^\infty(\Omega)}$. 它如何依赖于层数 $i$?
\end{MBExercise}

\begin{MBExercise}
\label{exr:ch09-04}
证明在尺寸为 $h$ 的规则网格上, 第~\ref{sec:ch09-condition-number-bounds} 节中的矩阵 $\mathbf A$ 满足 $\kappa(\mathbf A)=O(h^{-2})$.
\end{MBExercise}

\begin{MBExercise}
\label{exr:ch09-05}
令 $T=\{(x,y):x,y>0;\ x+y<1\}$ 是直角三角形, $V_T=\mathring H^1(T)$. 证明对任意次数为 $r$ 的多项式 $P$,
\[
\MathBookFormulaID{formula-ch09-ex05-volume-polynomial-inf-sup}
\sup_{v\in V_T}\frac{\int_TvP\,dx}{|v|_{H^1(T)}}
\geq c_r\|P\|_{L^2(T)},
\]
其中 $c_r$ 只依赖于次数 $r$. 提示: 只需让 $V_T$ 由固定次数的多项式组成; 使用有限维空间上范数的等价性, 并参见式~\MBUnresolvedReference{eq:ch12-5-2}{12.5.2}.
\end{MBExercise}

\begin{MBExercise}
\label{exr:ch09-06}
在 $r=0$, 即 $P$ 为常数的情形, 通过显式构造证明习题~\ref{exr:ch09-05} 的结果. 提示: 取 $v$ 为泡函数, 例如在 $T$ 内部为正、在 $\partial T$ 上为零的三次多项式.
\end{MBExercise}

\begin{MBExercise}
\label{exr:ch09-07}
令 $T=\{(x,y):x,y>0;\ x+y<1\}$ 是直角三角形, $e$ 是它的一条边. 令 $V_e$ 为在 $\partial T\setminus e$ 上为零, 且在 $T$ 上与次数为 $r$ 的多项式正交的函数空间. 证明对任意次数为 $r$ 的多项式 $P$,
\[
\MathBookFormulaID{formula-ch09-ex07-face-polynomial-inf-sup}
\sup_{v\in V_e}\frac{\int_evP\,ds}{|v|_{H^1(T)}}
\geq c_r\|P\|_{L^2(e)},
\]
其中 $c_r$ 只依赖于次数 $r$. 提示: 只需让 $V_e$ 由固定次数的多项式组成; 使用有限维空间上范数的等价性, 并参见式~\MBUnresolvedReference{eq:ch12-5-2}{12.5.2}.
\end{MBExercise}

\begin{MBExercise}
\label{exr:ch09-08}
在 $r=0$, 即 $P$ 为常数的情形, 通过显式构造证明习题~\ref{exr:ch09-07} 的结果. 提示: 考虑三次有限元, 在各边上取通常的 Lagrange 节点, 其余节点值取三角形上的积分. 构造一个在 $e$ 内部为正的 $V_e$ 中三次函数.
\end{MBExercise}

\begin{MBExercise}
\label{exr:ch09-09}
在 $d\geq3$ 维、$T$ 为单纯形的情形, 表述并证明习题~\ref{exr:ch09-05} 和~\ref{exr:ch09-07} 的结果.
\end{MBExercise}

\begin{MBExercise}
\label{exr:ch09-10}
设 $\alpha$ 和 $f$ 是固定次数 $r$ 的分片多项式, 证明定理~\ref{thm:ch09-residual-estimator-upper-bound}, 其中常数 $c$ 只依赖于 $r$ 和 $\mathcal T^h$ 的粗壮度. 提示: 使用习题~\ref{exr:ch09-05} 和~\ref{exr:ch09-07}.
\end{MBExercise}

\begin{MBExercise}
\label{exr:ch09-11}
证明若 $|h|_{W_\infty^1(\Omega)}$ 足够小, 且正则性估计~\eqref{eq:ch05-elliptic-regularity-h2} 成立, 则
\[
\MathBookFormulaID{formula-ch09-ex11-l2-estimate}
\|u-u_h\|_{L^2(\Omega)}\leq C\|\sqrt\alpha h^2\nabla^2u\|_{L^2(\Omega)}.
\]
提示: 参见第~\ref{sec:ch00-weighted-norm-estimates} 节.
\end{MBExercise}

\MathBookSourcePage{280}
\begin{MBExercise}
\label{exr:ch09-12}
证明当 $\alpha$ 和 $f$ 是次数为 $r$ 的分片多项式时, $h_T^{-n/2}E_T(u_h)$ 与 $E_T^\infty(u_h)$ 等价, 即
\[
\MathBookFormulaID{formula-ch09-ex12-estimator-equivalence}
\frac1{c_r}E_T^\infty(u_h)
\leq h_T^{-n/2}E_T(u_h)
\leq c_rE_T^\infty(u_h)
\qquad\forall T,
\]
其中 $c_r$ 只依赖于次数 $r$. 提示: 使用逆估计和 Hölder 不等式.
\end{MBExercise}

\begin{MBExercise}
\label{exr:ch09-13}
证明当 $\alpha$ 和 $f$ 是次数为 $r$ 的分片多项式时, $h_T^{n/2}E_T(u_h)$ 与 $E_T^1(u_h)$ 等价, 即
\[
\MathBookFormulaID{formula-ch09-ex13-estimator-equivalence}
\frac1{c_r}E_T^1(u_h)
\leq h_T^{n/2}E_T(u_h)
\leq c_rE_T^1(u_h)
\qquad\forall T,
\]
其中 $c_r$ 只依赖于次数 $r$. 提示: 使用逆估计和 Hölder 不等式.
\end{MBExercise}

\begin{MBExercise}
\label{exr:ch09-14}
证明式~\eqref{eq:ch09-polynomial-volume-inf-sup} 中必须限制 $P$ 所属的空间, 方法是证明
\[
\MathBookFormulaID{formula-ch09-ex14-unrestricted-inf-sup-zero}
\inf_{P\in L^2(T)}\sup_{v\in V_T}
\frac{\int_TvP\,dx}{\|P\|_{L^2(T)}|v|_{H^1(T)}}=0,
\]
其中 $T,V_T$ 与习题~\ref{exr:ch09-05} 相同. 提示: 令 $\mathcal T^h$ 是 $T$ 的一个剖分, 取 $P$ 在 $\mathcal T^h$ 的每个三角形上与常数正交. 用 $\mathcal T^h$ 上的分片常数逼近 $v$, 再令 $h\to0$.
\end{MBExercise}

\begin{MBExercise}
\label{exr:ch09-15}
回顾第~\ref{sec:ch08-main-theorem} 节和第~\ref{sec:ch08-reduction-to-weighted-estimates} 节中的 $\sigma_z$ 与 $\delta_z$. 证明
\[
\MathBookFormulaID{formula-ch09-ex15-weighted-delta-estimate}
\int_\Omega\sigma_z^{n+\lambda}|\nabla\delta_z(x)|^2\,dx
+h_z^{-2}\int_\Omega\sigma_z^{n+\lambda}|\delta_z(x)|^2\,dx
\leq Ch_z^{\lambda-2}.
\]
\end{MBExercise}

\begin{MBExercise}
\label{exr:ch09-16}
回顾第~\ref{sec:ch08-main-theorem} 节中的 $\sigma_z$. 证明
\[
\MathBookFormulaID{formula-ch09-ex16-inverse-weight-integral}
\int_\Omega\sigma_z^{-n-\lambda}\,dx\leq Ch_z^{-\lambda}.
\]
\end{MBExercise}

\begin{MBExercise}
\label{exr:ch09-17}
设 $e$ 是 $\mathcal T^j$ 的一个内部面. 证明在第~\ref{sec:ch09-convergent-adaptive-algorithm} 节的假设下,
\[
\MathBookFormulaID{formula-ch09-ex17-face-estimator-reduction}
E_e(u_h)^2
\leq C\sum_{T\in T_e}
\left(\|u_h-u_{h'}\|_{H^1(T)}^2+\operatorname{osc}(T)^2\right),
\]
其中正常数 $C$ 只依赖于系数 $\alpha$ 和 $\mathcal T^h$ 的形状正则性.
\end{MBExercise}

\begin{MBExercise}
\label{exr:ch09-18}
设 $D_1\subset D_2$ 是 $\mathbb R^n$ 中两个非退化单纯形. 证明
\[
\MathBookFormulaID{formula-ch09-ex18-polynomial-domain-norm}
\|p\|_{L^2(D_2)}\leq C\|p\|_{L^2(D_1)}
\qquad\forall p\in\mathcal P_k,
\]
其中正常数 $C$ 只依赖于 $k$, $D_1,D_2$ 的形状和比值 $|D_2|/|D_1|$. 用这一结果说明式~\eqref{eq:ch09-projected-residual-bubble-bound} 中第一个不等式和式~\eqref{eq:ch09-jump-bubble-identity} 中第二个不等式.
\end{MBExercise}

\begin{MBExercise}
\label{exr:ch09-19}
证明例~\ref{ex:ch09-refinement-without-energy-reduction} 中的断言.
\end{MBExercise}
